\section{Per Source Resolve Rate for Model and Loop Configurations}\label{app:per-source}

The main experiments report aggregate Resolve Rate and loop trace metrics over the 112 \lhb{} tasks. This appendix reports the base RQ1 setting before outer continuation, matching the \textit{w/o} RR column in Table~\ref{tab:rq1}. The three task sources represent distinct evidence regimes with different routing burden, state continuity, and regression obligation retention profiles. PR Sequences inherit source evidenced dependency DAGs from merged commit history, Course Labs inherit module level dependency structure from curriculum design, and Research Evolutions inherit coarser successor chains from citation graphs. Per source reporting separates these loop engineering patterns before aggregation.
Table~\ref{tab:per-source-rr} reports per source Resolve Rate for every model and loop configuration evaluated in RQ1. This split places the aggregate comparison in source context and relates each source to the routing and obligation retention burdens faced by each loop implementation. The aggregate column reproduces the corresponding \textit{w/o} number from Table~\ref{tab:rq1} and equals $\text{RR}_{\text{agg}} = \tfrac{1}{112}\big(29\cdot\text{RR}_{\text{PR}} + 57\cdot\text{RR}_{\text{Course}} + 26\cdot\text{RR}_{\text{Research}}\big)$ up to rounding.
Two findings hold uniformly across model and loop configurations and frame the interpretation of the per source columns.
\textbf{Finding 1: PR Sequences are uniformly the most difficult source.}
Across all 14 model and loop configurations, PR Sequence Resolve Rate reaches at most 3.45\%, and 9 of 14 configurations resolve no PR Sequence tasks. The PR Sequence pool is dominated by merge chains from actively maintained repositories, where each segment often depends on cross module invariants established several PRs earlier. When an early segment fails to preserve a prerequisite invariant, downstream progress becomes governed by obligation retention across dependent edits. This pattern is consistent with the RQ3 observation that regression events remain present across loop profiles, including runs with no recorded internal compaction.
\textbf{Finding 2: Course Labs account for most resolved tasks.}
For every model and loop configuration, Course Labs account for $\geq 70\%$ of the resolved task count. Course Labs have shallower dependency depth on average (median depth $3$ vs.\ $6$ for the pool) and a curriculum imposed module decomposition that often makes the human development order recoverable from task structure. Per source Resolve Rate on Course Labs varies most with model and loop configuration differences, and it tracks the aggregate Resolve Rate ranking most closely.
Research Evolutions fall between the two extremes. The moderate pool size ($26$ tasks) provides a more reliable per source estimate than earlier preview cuts, while the strongest model and loop configurations resolve a limited number of tasks and the weakest configurations resolve none. At this sample size, the per source spread is read together with the column level rate.

\begin{table}[tp]
  \centering
  \caption{Resolve Rate by task source before outer continuation.}
  \label{tab:per-source-rr}
  \small
  \setlength{\tabcolsep}{4pt}
  \renewcommand{\arraystretch}{1.08}
  \begin{tabular}{l l c c c c}
    \toprule
    \textbf{Loop} &  \textbf{Model}
      & \textbf{PR Seq.}\,(n{=}29) & \textbf{Course}\,(n{=}57) & \textbf{Research}\,(n{=}26)
      & \textbf{All}\,(n{=}112) \\
    \midrule
    \multicolumn{6}{@{}p{0.98\linewidth}@{}}{\emph{Model sweep under the Claude Code loop}} \\
    Claude Code & Opus-4.7        & 3.45\%  & 24.56\% & 15.38\% & 16.96\% \\
    Claude Code & GPT-5.5         & 0.00\%  & 21.05\% & 11.54\% & 13.39\% \\
    Claude Code & GLM-5.1         & 0.00\%  & 21.05\% & 11.54\% & 13.39\% \\
    Claude Code & DeepSeek-V4P    & 0.00\%  & 19.30\% & 7.69\%  & 11.61\% \\
    Claude Code & Gemini-3.1-Pro  & 0.00\%  & 14.04\% & 7.69\%  & 8.93\%  \\
    Claude Code & Qwen3.6-Plus    & 0.00\%  & 10.53\% & 3.85\%  & 6.25\%  \\
    Claude Code & Kimi-2.6        & 0.00\%  & 5.26\%  & 3.85\%  & 3.57\%  \\
    Claude Code & Grok-4.1-FR     & 0.00\%  & 5.26\%  & 0.00\%  & 2.68\%  \\
    \midrule
    \multicolumn{6}{@{}p{0.98\linewidth}@{}}{\emph{Loop sweep under a fixed \texttt{gpt-5.4} model}} \\
    Codex          & gpt-5.4 & 0.00\%  & 21.05\% & 11.54\% & 13.39\% \\
    Claude Code    & gpt-5.4 & 0.00\%  & 19.30\% & 11.54\% & 12.50\% \\
    GitHub Copilot & gpt-5.4 & 0.00\%  & 17.54\% & 7.69\%  & 10.71\% \\
    OpenHands      & gpt-5.4 & 0.00\%  & 10.53\% & 3.85\%  & 6.25\%  \\
    SWE-agent      & gpt-5.4 & 0.00\%  & 8.77\%  & 3.85\%  & 5.36\%  \\
    mini-swe-agent & gpt-5.4 & 0.00\%  & 8.77\%  & 0.00\%  & 4.46\%  \\
    \bottomrule
  \end{tabular}
\end{table}

These patterns shape the interpretation of aggregate \lhb{} Resolve Rate. A single aggregate comparison combines Course Labs variation with the limited spread in PR Sequences, and the loop sweep gap in Table~\ref{tab:rq1} is driven largely by Course Labs. PR Sequence gains are clearer in the per source columns than in aggregate deltas. Doubling PR Sequence RR moves the aggregate by less than one percentage point while materially changing the loop engineering interpretation of routing burden and obligation retention.

\FloatBarrier
